\chapter{Machine Settings}
\label{ch:settings}

The Settings page holds values that belong to the machine rather than to any
recipe: lighting, axis speeds, clamp voltage, and the calibration entries.
They apply whichever configuration is loaded.

\section{Getting there}

From the configuration list, press \keycap{SAVE}. Press \keycap{ESC/DEL} to
leave.

\begin{wbnote}
\keycap{SAVE} opens the Settings page from the configuration list, and saves
the Settings page once you are on it. It does two different jobs depending on
where you are.
\end{wbnote}

\section{The settings}

% The Settings page -- machine-wide values, shared by every configuration.
%
% DERIVED FROM FIRMWARE. Refresh against:
%   Firmware/Core/Src/menu.cpp                  (row labels and order)
%   Firmware/Core/Src/user_interface_module.cpp (kSettingLimits)
%   Firmware/Core/Inc/configuration.h           (min/max/step/default macros)
%   Firmware/Core/Inc/machine_settings.hpp      (field semantics)

\begin{center}\small
\begin{tabular}{@{}L{8mm}L{32mm}L{12mm}R{16mm}R{16mm}R{13mm}R{16mm}@{}}
\toprule
\textbf{Row} & \textbf{Label} & \textbf{Unit} & \textbf{Minimum} &
\textbf{Maximum} & \textbf{Step} & \textbf{Default} \\
\midrule
1 & \menuitem{Clamp Volt}       & V    & 8.0 & 12.0  & 0.1 & 8.0 \\
2 & \menuitem{Area Light}       & \%   & 0   & 100   & 1   & 100 \\
3 & \menuitem{Spotlight}        & \%   & 0   & 100   & 1   & 100 \\
4 & \menuitem{Up Speed}         & mm/s & 0.1 & 10.0  & 0.1 & 7.5 \\
5 & \menuitem{Down Speed}       & mm/s & 0.1 & 10.0  & 0.1 & 5.0 \\
6 & \menuitem{Setup Force}      & g    & 0.0 & 100.0 & 0.5 & 30.0 \\
7 & \menuitem{Spot On}          & \multicolumn{5}{@{}l@{}}{on/off --- \keycap{ENTER} toggles} \\
8 & \menuitem{Start Tach. Cal.} & \multicolumn{5}{@{}l@{}}{action --- \keycap{ENTER} runs it} \\
\bottomrule
\end{tabular}
\end{center}

\begin{description}[leftmargin=!,labelwidth=30mm,style=nextline,font=\normalfont\ttfamily]
\item[Clamp Volt] Holding voltage applied to the wire clamp solenoid. Raise it
  if the clamp slips under wire tension; keep it as low as will hold, since the
  solenoid runs cooler.
\item[Area Light] Brightness of the work area illumination, as a percentage.
  The \keycap{LIGHT} button switches it on and off; this sets how bright it is
  when on.
\item[Spotlight] Brightness of the spotlight, as a percentage.
\item[Up Speed] Maximum speed the Z axis is allowed to rise at. Higher is
  faster per cycle; too high can shake the loop.
\item[Down Speed] Maximum speed the Z axis is allowed to descend at. Keep this
  moderate --- the head must be able to stop on contact.
\item[Spot On] Whether the spotlight is enabled at all. There is no panel
  button for the spotlight; this row is the only switch.
\item[Setup Force] Force the machine holds during a calibration run started
  with \keycap{SETUP}. This is the value your gauge reading is compared
  against.
\item[Start Tach. Cal.] Runs tachometer calibration. See
  Chapter~\ref{ch:calibration}.
\end{description}

\begin{wbnote}
Changes on this page take effect the moment you make them --- the light
brightens, the clamp voltage changes --- but they are \emph{not} remembered
until you press \keycap{SAVE}. Leaving with \keycap{ESC/DEL} keeps the live
values for this session but restores the saved ones at the next power-up.
\end{wbnote}

\begin{wbnote}
Two more values are stored with these settings but never appear on the page,
because the machine writes them for you: the \textbf{force offset} from a
\keycap{SETUP} run, and the \textbf{tachometer offset} from a calibration run.
Both are saved as soon as they are produced.
\end{wbnote}

\derivedfrom{\texttt{menu.cpp}, \texttt{user\_interface\_module.cpp},
             \texttt{machine\_settings.hpp}, \texttt{configuration.h}}


\section{Choosing axis speeds}

\menuitem{Up Speed} and \menuitem{Down Speed} cap how fast the bond head
moves.

\begin{description}[leftmargin=!,labelwidth=30mm,style=nextline,font=\normalfont\ttfamily]
\item[Up Speed] Mostly affects cycle time. Raise it until the loop starts to
  show signs of being disturbed by the acceleration, then back off.
\item[Down Speed] Affects bond quality. The head must be able to stop cleanly
  on contact; too fast and it overshoots into the pad before the force loop
  catches it. Slower is safer. The default of 5\,mm/s is a reasonable starting
  point.
\end{description}

\begin{wbwarning}
Raising \menuitem{Down Speed} substantially increases the impact energy at
touchdown and can crater pads. Change it in small steps and verify bond quality
after each change.
\end{wbwarning}

\section{Clamp voltage}

\menuitem{Clamp Volt} sets how hard the wire clamp holds. Use the lowest
voltage that reliably holds the wire through the tear: the solenoid runs cooler
and lasts longer. Raise it if the wire slips during the tear.
